\clearpage
\chapter*{Đính chính cho Bài trình bày số 182}
\addcontentsline{toc}{chapter}{Đính chính cho Bài trình bày số 182}
\setcounter{footnote}{1}
\sid{e182}
\markboth{Đính chính 182}{Đính chính 182}
\src{298}

\noindent\textbf{2. Hình học hình thức và hình học đại số.}

\begin{center}
[Hội thảo Bourbaki, t. 11, 1958/59, số 182, 28 tr.]
\end{center}

\noindent\textbf{Trang 8, Định lý 6.} --- Giả thiết «$f$ tựa xạ ảnh» có thể
được thay bằng giả thiết yếu hơn «$f$ tách», nhờ kết quả sau đây (xem SGA VIII,
6.2): mọi cấu xạ $f:X\longrightarrow Y$ tựa hữu hạn và tách đều là tựa xạ ảnh.

\src{299}
\noindent\textbf{Trang 12, và trang 14, Nhận xét 1.} --- Xin lưu ý rằng J.-P.
Serre đã xây dựng một đa tạp xạ ảnh không kỳ dị, chiều 3, trên một trường đóng
đại số $k$ có đặc số $p>0$, không thu được bằng phép rút gọn từ một lược đồ
proper trên một vành địa phương nguyên có trường thặng dư $k$ và trường phân
thức có đặc số không.\footnote{Serre (Jean-Pierre). --- Exemples de variétés
projectives en caractéristique $p$ non relevables en caractéristique zéro,
\emph{Proc. Nat. Acad. Sc. U.S.A.}, t. 47, 1961, tr. 108--109.} Mumford được
cho là đã tìm được một ví dụ tương tự với một mặt xạ ảnh không kỳ dị.

\noindent\textbf{Trang 15, các Nhận xét 2 và 3.} --- Hiện nay tôi ít lạc quan
hơn về những kết quả được phỏng đoán ở đây. Tuy nhiên, bài toán về các biến
thiên cấu trúc của không gian xạ ảnh, được nêu ở cuối Nhận xét 3, đã được
Hironaka giải quyết theo hướng khẳng định, và bài toán tương tự đối với các đa
tạp abel đã được Koizumi giải quyết.

\noindent\textbf{Trang 16, dòng 3.} --- Xin lưu ý rằng Mumford gần đây đã xây
dựng các lược đồ môđuli cho các đường cong giống $g$ (xem SHMI). Ngoài ra,
Định lý 10 chứng minh rằng các «lược đồ Jacobi mức $n$» trong lý thuyết môđuli
là không kỳ dị (và thậm chí trơn trên $\mathbf Z$).

\noindent\textbf{Trang 23, dòng 12, dòng đầu của Hệ quả 5.} --- Bổ sung:
«$X$ hoặc $Y$ proper trên $k$».

Nội dung cốt lõi của các §§ 1 đến 5 nằm trong phần EGA III đã xuất bản, còn
nội dung của các §§ 6 và 7 nằm trong SGA III. Để nghiên cứu nhóm cơ bản, xem
SGA V, IX, X, XI, cũng như SGA 1962 (các Bài trình bày X, XII và XIII), về các
định lý kiểu Lefschetz và nhiều câu hỏi mở. Chỉ riêng lý thuyết về các phủ phân
nhánh thuần (xem Định lý 14) là chưa được trình bày thành một bản hoàn chỉnh.
Hệ quả của Định lý 14, cho phép xác định hoàn toàn các phủ Galois có cấp nguyên
tố cùng nhau với đặc số của một đường cong đại số trên một trường đóng đại số,
đã được sử dụng một cách thiết yếu trong ba trường hợp:

1\textsuperscript{o} trong chứng minh của Igusa về bất đẳng thức Picard cho
các mặt xạ ảnh không kỳ dị trong đặc số tùy ý;

2\textsuperscript{o} trong việc nghiên cứu (do Ogg và Šafarevič phát triển độc
lập) nhóm các phân thớ chính thuần nhất dưới một đa tạp abel được xác định trên
một trường hàm một biến, trong đặc số tùy ý;

3\textsuperscript{o} trong chứng minh gần đây của Artin về một số định lý then
chốt của «đối đồng điều Weil» của các đa tạp đại số.

\par
